\documentclass[11pt]{article}
\usepackage{amsmath,amssymb,amsthm,hyperref}
\newtheorem{theorem}{Theorem}\newtheorem{lemma}[theorem]{Lemma}\newtheorem{proposition}[theorem]{Proposition}
\newcommand{\R}{\mathsf R}\newcommand{\Col}{\mathsf{Col}}\newcommand{\Vect}{\mathbf{Vect}_k}
\title{Typed polynomial recursive feedback with retained native grades}
\date{16 November 2026}
\begin{document}\maketitle
\paragraph{Status.} The written proof has been independently reviewed at the declared scope; no formal verification is claimed. This constructs recursive execution feedback for actual finite-occurrence polynomial content, rather than replacing that content by one-hole writer output. The positive result is indexed by a specified countable augmented-rig extension and an independently presented class of polynomial contentads. Its global content grade retains internal recursive-edge labels, including unused shapes. Finite terminating execution effects exist without almost-sure termination; their missing mass is not identified with divergence. Instantaneous interconnection, arbitrary shared branching effects and the full v0.2 completion remain separate.

\section{Source reconstruction and the two material changes}
The starting source is \emph{General Independent Interaction Criterion}, repository path
\path{projects/contextual-composition/interaction/independent-criterion/General_Independent_Interaction_Criterion.tex}, blob
\texttt{1315e4b4fca69e1899498ae831f4fe41b60b4d52}.
The required passages are its augmented-family native arrows and dependent Fubini construction, L1--L6, and its constant ordered-polynomial content example. The later all-coupling span bridge is not assumed: it changes conditional bind oplaxly and does not supply the polynomial crossing used here.

Two changes are declared in advance. Contexts admit countable raw witness sets and subprobability rather than only finite normalized effects. Content after closing a recursive wire $U$ is not the original collection $C$, but a collection $\R_U C$ of finite trees whose internal edges retain their actual $U$ labels. Both changes are essential: a finite-support effect need not contain all terminating executions, and removing the internal edge labels loses the global-grade inverse needed for nested feedback. This operation is not a trace on a fixed content grade.

\section{An actual countable extension of an original augmented-rig source}
Fix a field $k$. The finite normalized source uses the symmetric rig
\[
 \mathcal V_0=(\mathbb R_{\geq0})_{\rm disc}\times\mathbf{FinVect}_k,
\]
with addition/direct sum and multiplication/tensor. Its augmented family has a finite raw set $W$, masses $p_w$ summing to one, vector spaces $V_w$, and an augmentation $\bigoplus_wV_w\to k$. Since direct sum is a coproduct, the last datum is equivalently actual arrows $\epsilon_w:V_w\to k$. A native map $r:W\to W'$ has exact mass pushforward and actual maps
\[
 \alpha_{w'}:\bigoplus_{r(w)=w'}V_w\longrightarrow V'_{w'},
 \qquad \epsilon'_{w'}\alpha_{w'}=(\epsilon_w)_{r(w)=w'}.
\]
These maps can be nonfaithful under the raw-carrier interpretation. They are equivalently their individual restrictions $V_w\to V'_{r(w)}$, but not merely the raw function.

For the target use
\[
 \mathcal V_\omega=([0,\infty],\leq)\times\Vect.
\]
The scalar category has one arrow precisely when the inequality holds. Addition includes countable nonnegative sums, multiplication uses $0\cdot\infty=0$, and the vector factor uses countable direct sums and ordinary tensor products. Contexts have countable raw sets and global augmentation to $(1,k)$. Thus their scalar datum is exactly $\sum_wp_w\leq1$, and their vector datum is the same family of actual augmentations. Native maps satisfy
\[
 \sum_{r(w)=w'}p_w\leq p'_{w'}
\]
and the displayed vector augmentation equation. In particular the extension permits mass-subpreserving cells; it does not describe them as old exact cells.

\begin{lemma}[Countable crossing and finite-source recovery]\label{lem:crossing}
The augmented context unit, dependent accumulation, finite occurrence tensor and native Fubini of the original source extend to these countable augmented families. For any finitary cartesian polynomial occurrence monad and independent content satisfying L1--L6 over countable-fibre maps, they give the same native crossing with countable context fibres. The finite normalized $\mathcal V_0$ source embeds locally fully faithfully and strongly for its original dependent and occurrence operations.
\end{lemma}
\begin{proof}
Dependent witness sets are countable sums of countable sets, and a finite occurrence tensor has a finite product of countable sets. Their weights and vector coefficients are respectively products and finite tensors. Countable nonnegative distributivity and Fubini prove the scalar augmentation inequalities. On vectors, tensor preserves direct sums, giving the same actual aggregation formula as the original dependent Fubini construction, now with countable fibre sums. To verify its identities, composition and reassociation equations, restrict both candidate linear maps to each direct-sum summand. The restrictions are the same finite tensor of the original arrows and augmentations, with the prescribed addressed permutation; symmetric monoidal coherence proves equality. Direct-sum injections jointly determine a map, and the scalar factor is thin, so these are full native equalities, not only carrier identities.

Countable-fibre maps are closed under pullback and composition. A finitary polynomial functor preserves that class because the fibre over a fixed shape and labelling is a finite product of the original fibres. The original display construction therefore produces an admitted context. Its content comparisons follow from the same independently supplied cartesian/opcartesian universal properties and L1--L6. The context-side structural diagrams have just been checked on actual countable aggregation maps, so no infinite distributivity axiom is merely inferred from the old finite theorem.

For two old normalized finite families, the target scalar inequalities all have equal total one. Hence each fibre inequality is equality: the nonnegative deficits sum to zero. The vector maps and their augmentation equations are unchanged under the fully faithful inclusion of finite-dimensional spaces in all vector spaces. This proves local fullness and faithfulness between old context objects; finite sums, finite tensors and their actual structural maps are preserved, proving strong recovery of the old operations. The unchanged content input and body equations then give the same local recovery in the mixed comma presentation.
\end{proof}
The statement is about this scalar/vector original rig, not all original rigs. It uses ordinary set-theoretic countable choice. Every augmented scalar is finite because its total is at most one; infinite scalar objects are included only to make the ambient countable semiring closed under its declared operations.

\section{Independently supplied nontrivial polynomial content}
Let $\mathcal D$ be a small or locally small symmetric monoidal category and let $\Col(\mathcal D)$ be ordered collections of families: an object $C$ has a set $S_n$ of shapes of each finite arity $n$ and an actual coefficient $D_c\in\mathcal D$ for each shape. A native map preserves arity, sends each shape $c$ to a shape $a(c)$, and includes an actual arrow $D_c\to D'_{a(c)}$. Its extension is
\[
 F_C(B)=\coprod_{n\geq0}S_n\times B^n,
\]
forgetting the coefficient but not the shape. Unused shapes and coefficient arrows remain part of the global grade.

Use the constant fibration $\mathbf{Set}\times\Col(\mathcal D)\to\mathbf{Set}$, constant pure transport, and occurrence map $F_C(B)\to\operatorname{List}(B)$ forgetting only the shape. The unit collection has one unary shape with coefficient the tensor unit. Accumulation $C\diamond D$ has shapes of $F_D(F_C(B))$: an outer $D$ shape with an ordered list of inner $C$ shapes. Its coefficient is the outer coefficient tensored with all the inner coefficients; the body is polynomial substitution.

\begin{lemma}
These data satisfy the original independent content conditions L1--L6 for all set maps. Native maps fix the ordered positions and may merge shapes of the same arity; arbitrary permutations or deletion of occurrence positions are not included. This is the precise ordered-content restriction in the original source example. With $\mathcal D=\mathbf{FinVect}_k/k$ they include actual nonfaithful native content coefficient maps, and the unit grade has only its identity endomorphism.
\end{lemma}
\begin{proof}
The original ordered-polynomial example checks the shape and extension part. On coefficients, nested substitution lists each labelled node once. Reassociation uses the prescribed node-address bijection and the actual symmetry/associator maps in $\mathcal D$. All pentagon and unit paths induce the same addressed permutation, so symmetric monoidal coherence proves their equality; naturality follows by tensoring the supplied coefficient arrows. This verifies the one-sided contentad, not just a collection of object formulas. Cartesian arrows in the constant fibration are exactly arrows with invertible grade component. Unit and accumulation preserve them. Display is the pullback that retains one shape and relabels its finite ordered positions; Beck--Chevalley is identity on grades. The occurrence equations are singleton and list flattening, unchanged by coefficients. Finally the tensor unit in $\mathbf{FinVect}_k/k$ is terminal, while augmented vector spaces can have distinct arrows with identical shape maps. Thus unit rigidity and substantive nonfaithful grades coexist.
\end{proof}

\section{The typed global feedback grade}
For a set $U$, define $\R_U C$ as follows. A shape is a nonempty finite rooted planar tree, with each node labelled by a shape of $C$ having exactly its number of ordered outgoing slots. A slot is either an external leaf or an internal edge to another node; each internal edge has a label in $U$. The arity is the number of external leaves, in their planar order. The coefficient of this shape is the finite tensor of all node coefficients, using a chosen planar ordering. All such shapes are retained, regardless of whether a process body can ever return them.

No zero-node identity tree is included. The occurrence content still has its separate unary unit collection. Including an arbitrary zero-node block in $\R_U C$ would destroy the unique block decomposition below.

\begin{theorem}[Initial algebra and coherent nesting]\label{thm:grade}
In the native category of coefficient-decorated ordered collections, $\R_U C$ is the initial algebra of
\[
 X\longmapsto C\circ(\mathsf I+U\cdot X).
\]
Here the sum is tagged family sum, $\mathsf I$ is the unary unit collection, and $\circ$ is ordinary polynomial substitution with the left factor outermost (the reverse notational order to the source execution symbol $\diamond$). Its coefficients are obtained by the corresponding finite tensor of node coefficients. It is functorial in $U$ and all actual native maps of $C$. There are canonical native isomorphisms
\[
 \R_\varnothing C\cong C,
 \qquad
 \R_U(\R_V C)\cong\R_{U+V}C.
 \tag{G}
\]
Their unit, associativity and parameter-naturality coherence holds on shapes and actual coefficient arrows.
\end{theorem}
\begin{proof}
Root decomposition gives the algebra structure. Recursion on finite node count gives its unique shape map to any native algebra of the displayed polynomial endofunctor. At a node, tensor the already determined child coefficient arrows and compose with that algebra's actual root coefficient arrow. This recursively determines a native algebra morphism; any such morphism must have exactly these shape and coefficient components by induction on node count. This proves initiality on actual native arrows, not just on shape sets. A map of $U$ relabels internal edges; a native map of $C$ acts on every node and tensors the actual node coefficient arrows.

With no internal edge labels available, every tree is a single corolla, giving the first isomorphism. For the second, flatten a tree whose nodes are $V$-edge-labelled $C$ trees, labelling the edges joining distinct blocks by $U$ and the edges within a block by $V$. Conversely cut every $U$-labelled edge of a $U+V$-tree. Its connected components are nonempty $V$-trees, and their contracted outer tree has $U$-labelled edges. These operations are mutually inverse, including nullary nodes and zero external arity. Coefficients are related by the actual monoidal isomorphism for the corresponding labelled-node bijection. For three loop sets, both flattening routes give the identical tree with its original three colours on edges; both coefficient routes implement the same node permutation. This proves coherence by symmetric monoidal coherence, not by forgetting coefficients. Relabelling or mapping the original shapes commutes with the construction, proving naturality.
\end{proof}

Initiality here is the elementary finite-tree algebra property. It is not claimed as a universal property of the whole mixed feedback extension. In particular no theorem says that a functor out of the original native process category automatically interprets these new grades or preserves the countable effects.

\section{Recursive effect and body construction}
Take a native mixed process
\[
 f:A+U\longrightarrow F_C(B+U)
\]
in the countable context/polynomial-content presentation above. Thus each input $x$ has countable raw witnesses $W_x$, subprobability weights $p_x(w)$, augmented vector coefficients $V_{x,w}\to k$, and a body returning a shape $c(w)$ and an ordered finite list in $B+U$.

A terminating execution witness rooted at $a\in A$ is a finite tree constructed by choosing a witness at its root, expanding each returned $U$ slot with a witness at that specified input, and leaving every returned $B$ slot as an external leaf. No $U$ slot may be left unexpanded. All raw zero-weight choices are retained. Write $\mathsf T_a(f)$ for this set. Its weight and coefficient are
\[
 p(t)=\prod_{v\in\operatorname{nodes}(t)}p_{x_v}(w_v),
 \qquad V_t=\bigotimes_v V_{x_v,w_v},
\]
with the tensor augmentation. Its content body retains the $C$ shape at every node and the actual $U$ value on each internal edge, and labels its external leaves by the returned $B$ values. It therefore lands in $F_{\R_U C}(B)$.

\begin{theorem}[Total finite-execution feedback]\label{thm:feedback}
The preceding data define a native mixed process
\[
 f^{\dagger U}:A\longrightarrow F_{\R_U C}(B)
\]
for every small loop set $U$, with no almost-sure termination or finite-support assumption on the set of terminating executions. Its raw witnesses are countable at each input, and their total weight is at most one. The construction is functorial on native process cells whose endpoint maps preserve the external/recursive partition and use the same map on recursive inputs and outputs. Such a cell includes an actual row comparison $r_u:W_u\to W'_{h(u)}$ with its effect inequality and full body square at EVERY recursive input $u$; a bare state function $h$ is insufficient. This is a local native simulation hypothesis, not arbitrary loop-map dinaturality. Exact primitive mass pushforward gives exact pushforward on terminating tree weights; subpreserving primitive comparisons give subpreserving tree comparisons. Actual context aggregation maps and the entire global content map are retained.
\end{theorem}
\begin{proof}
At each fixed finite height there are countably many trees: choose one of countably many root witnesses and then one of a finite product of countable child-tree sets. The union over finite heights remains countable. This proof is local at each input and does not require $U$ itself to be countable.

Put $q_x^{(0)}=0$, and recursively
\[
 q_x^{(n+1)}=\sum_{w\in W_x}p_x(w)
       \prod_{i:\,y_i(w)\in U}q_{y_i(w)}^{(n)}.
\]
The empty product is one. Induction identifies this with the mass of trees of height at most $n+1$, and gives $0\leq q_x^{(n)}\leq1$. Hence the mass of all finite trees is their increasing limit and is at most one. The direct sum of the actual tensor augmentations supplies $\bigoplus_tV_t\to k$, so the effect bound and vector arrow give exactly the required augmented context. Countable sums and finite products commute with increasing limits, so $q=\sup_nq^{(n)}$ is the least fixed point of the displayed polynomial subprobability operator. This establishes only the scalar least-fixed-point property; the native witnesses and coefficients were constructed separately.

Let a cell have witness maps $r_x$, actual augmented coefficient maps, a native shape map $a:C\to C'$, and recursive state map $h:U\to U'$. The full body square says that node arities agree, external values are mapped by their specified endpoint map, and recursive outputs are mapped by $h$. Consequently nodewise application of $r_x$ defines a map of terminating trees with precisely the same finite addressed-node set. Its content map is the globally defined $\R_h(a)$, including every unused tree shape and coefficient.

For a fixed target tree, induction on its node count proves the effect inequality. Sum over source root witnesses mapping to the target root; for each such root, the child fibres are independent products. Apply the child induction first, then the primitive root pushforward inequality. Countable nonnegative Fubini justifies rearranging these sums. If all primitive equations are exact, every step is equality. A target zero-weight witness need not have a source preimage; this is consistent with exact pushforward and does not delete that target witness.

For each source tree, tensor the actual pointwise augmented coefficient arrows along its node-address bijection. Aggregate these maps over the full source fibre by the universal property of the vector direct sum. Their augmentations commute because every node arrow preserves its augmentation. These are exactly the native fibre aggregation arrows. Identity and composition follow nodewise; tensor functoriality and direct-sum uniqueness identify the full aggregated maps. The content map obeys the same laws by Theorem~\ref{thm:grade}. Thus the assertion is about genuine native cells, not just inequalities of total termination probabilities.
\end{proof}

\section{Unfolding uses the original multi-occurrence crossing}
The tree construction is not merely an operation placed beside the old occurrence substitution. Let $\bar f:A+U\to F_{\R_U C}(B)$ denote the same terminating-tree construction allowing its root to be any input of $f$, including a recursive input. Form the independently available collection
\[
 D=\mathsf I+U\cdot\R_U C.
\]
Define one continuation process $G:B+U\to F_D(B)$ as follows. At $b\in B$, its context is the unit context and its body is the unary unit shape labelled by $b$. At $u\in U$, its context is the context of $\bar f$ at $u$, and its body is the corresponding tree shape in the $u$-tagged summand of $D$. The complete global grade is $D$ in every case; no unused unit or tagged tree shape is removed. This is a legitimate tagged-grade case construction, not an assertion that arbitrary case-dependent global grades are equal.

\begin{theorem}[Native polynomial unfolding]\label{thm:unfolding}
Compose $f$ with $G$ using the original independent occurrence crossing, extended by Lemma~\ref{lem:crossing}. Its global content collection is $C\circ D$, with the first step's $C$ shape outermost. Root decomposition supplies a full native content isomorphism
\[
 C\circ(\mathsf I+U\cdot\R_U C)\cong\R_U C
\]
under which that mixed composite is natively isomorphic to $\bar f$. The isomorphism is natural under the local simulation cells of Theorem~\ref{thm:feedback} and retains all actual context and content coefficient maps.
\end{theorem}
\begin{proof}
The original crossing selects one continuation witness at each ordered occurrence returned by the root witness. At an external $B$ occurrence that witness is the unit; at a recursive $U$ occurrence it is a terminating tree rooted at the actual returned state. These data are uniquely the root and ordered children of one terminating tree. Thus the context raw map is a bijection, including all zero-weight witnesses. The old dependent/Fubini coefficient is root coefficient tensor the finite occurrence tensor of continuation coefficients; the tree coefficient is the same addressed tensor with the prescribed reassociation and permutation. The effect is the same product, and the augmentation equations are the existing augmented tensor equations. Root decomposition identifies the entire global content collection, including its unused shapes and actual node coefficients, not only the returned content bodies. The actual native body square follows from the same root/child decomposition. Applying a native simulation to every root and child commutes with this decomposition; tensor functoriality and the already proved countable aggregation naturality give the full native naturality equation.
\end{proof}
This supplies the explicit feedback--occurrence relationship at the declared source class. Its continuation children are independent conditional on the retained root witness. Replacing that tensor by an unspecified shared joint is not licensed by this theorem.

\section{Nesting, zero feedback and the obstruction to ordinary trace laws}
For $f:A+(U+V)\to F_C(B+(U+V))$, first close $V$ while treating $A+U$ and $B+U$ as external, then close $U$.

\begin{theorem}[Native nested-feedback law]
There is a canonical native isomorphism, with the global grade comparison from (G),
\[
 (f^{\dagger V})^{\dagger U}\cong f^{\dagger(U+V)}.
\]
It is natural in native cells whose recursive map has the form $h_U+h_V$, preserving the staged $U/V$ partition, and coherent for three or more disjoint recursive interfaces with colour-preserving maps. A state map mixing the two blocks is covered by the single-loop comparison theorem only after simultaneous closure; staged naturality under such a map is not asserted. Closing the empty interface recovers the original process and its full global grade up to the canonical corolla identification.
\end{theorem}
\begin{proof}
An iterated witness is an outer $U$ tree whose nodes are terminating $V$ trees. Flatten it. Conversely cut every internal edge labelled in $U$ in a simultaneous terminating tree; each component is the unique maximal $V$ subtree, and the contracted tree is the outer $U$ execution. This is a bijection of all raw witnesses, including zero-weight ones. Products of primitive weights agree, since each node occurs once. The actual context tensor coefficient maps are the coherent monoidal isomorphisms for this node bijection. The full content grade, not merely its used part, has exactly the isomorphism (G). Body and augmentation squares therefore hold. Triple coherence and colour-preserving native-cell naturality follow from the identical coloured-node address calculation and symmetric monoidal coherence. With empty $U$, every execution consists of one root witness and no recursive child; its raw context and coefficient are unchanged, and $\R_\varnothing C=C$ globally.
\end{proof}

The edge labels are essential to the global-grade conclusion. If all internal edge labels are forgotten and both loop sets are replaced by an untyped tree star, nested grade flattening is generally many-to-one: a two-node unary tree can be treated as one two-node block or as two singleton blocks. The typed decomposition recovers the block cut from its $U/V$ colours. It does not recover the cut from the untyped shape alone.

Ordinary trace or Conway axioms are not asserted. Even with a single unary generator and singleton loop set, $\R_U C$ has one unused unary-chain shape at every positive length. A particular yanking-style execution may use only one chain length, but that does not make its entire grade isomorphic to the original singleton collection. Deleting unused chains would change the native content interface. More general sliding or substitution laws must therefore state their actual typed grade maps and preservation strength rather than importing an ungraded trace equation.

\section{One retained scope gives genuinely shared branching recurrence}
Let $\Lambda$ be a countable latent-scope set, with seed weights $\mu_a(\ell)\geq0$ satisfying $\sum_\ell\mu_a(\ell)\leq1$ and actual augmented seed coefficients $L_{a,\ell}\to k$ at external input $a$. Supply conditional primitive rows indexed by $(a,\ell)$ and $(u,\ell)$, with witnesses $w$, subprobability weights $p_{x,\ell}(w)$, augmented coefficients $E_{x,\ell,w}$ and polynomial bodies in $B+U$, all using the same complete global content collection $C$. No conditional law is inferred from a marginal; these are independent input data.

Build a single process with recursive interface $U\times\Lambda$. Its external row at $a$ first chooses $(\ell,w)$ with weight $\mu_a(\ell)p_{a,\ell}(w)$ and coefficient $L_{a,\ell}\otimes E_{a,\ell,w}$. Every returned recursive value $u$ is relabelled $(u,\ell)$. A recursive row at $(u,\ell)$ uses only the conditional witness weight $p_{u,\ell}(w)$ and coefficient $E_{u,\ell,w}$; it does not sample another seed. All its recursive outputs again retain that same $\ell$.

\begin{proposition}[Retained-scope polynomial recurrence]
The preceding process is in the target context/content doctrine. Its terminating-tree feedback has weights and actual context coefficients
\[
 \mu_a(\ell)\prod_{v}p_{x_v,\ell}(w_v),
 \qquad
 L_{a,\ell}\otimes\bigotimes_v E_{x_v,\ell,w_v},
\]
where the conditional root row is included in the node product. Every reachable internal edge carries the same latent value $\ell$. The seed coefficient occurs exactly once, not once per branch or node. The global content grade remains the complete $\R_{U\times\Lambda} C$, including unused shapes with mixed latent labels. Native seed maps using a single common function $h:\Lambda\to\Lambda'$ across all external inputs, one common recursive map $h_U$, one common external-output map, and one common full global content map $a:C\to C'$ (including all coefficient arrows), together with conditional-row native simulations over those same maps, induce native comparisons through this feedback. Input-dependent seed functions $h_a$ are not covered unless enough input data is also retained in the recursive state to define one state map.
\end{proposition}
\begin{proof}
The external row is the original dependent accumulation of the seed context and the conditional root context. Its augmentation and native maps are given by countable dependent Fubini. Every recursive row is already supplied. Induction on tree nodes shows that the latent label never changes; expanding the weight and coefficient definitions then exhibits exactly one seed factor, at the external root. The rest is Theorem~\ref{thm:feedback}. For a seed function $h:\Lambda\to\Lambda'$, the required data include its native seed pushforward/augmentation equation and, for every $\ell$, an actual conditional-row comparison to the row at $h(\ell)$ with its full body square over the common $h_U$, external-output map and global content map $a$. Independently chosen content maps $a_\ell$ do not suffice to assemble one native process cell. These derive the external-row comparison by dependent Fubini and the recursive-row comparisons directly. The recursive state map is $(u,\ell)\mapsto(h_U(u),h(\ell))$, so the feedback theorem applies. A seed marginal alone supplies none of these conditional comparisons.
\end{proof}

For a concrete discriminator, sample a fair bit $\ell$ once at the external root and return a binary shape with two recursive children $(u,\ell)$. At each recursive child, with probability $0<p<1$ return a unary shape with external output $\ell$; with probability $1-p$ return another binary shape with both recursive outputs $(u,\ell)$. This permits unbounded branching recurrence. On the event that the first two children both terminate, of weight $p^2$, their two external bits agree with probability one. Independently resampling the bit at those children instead gives unequal bits with weight $p^2/2$. Thus the retained-scope construction is not the independent-mixture substitute. Nontrivial $L_\ell$ and $E_{u,\ell,w}$ make the once-only seed coefficient equally explicit at the native interface.

This is a fixed shared environment and correlated sibling execution. It is not arbitrary mutable shared state or every possible interfering coupling. The children are independent conditional on the retained seed, and the theorem preserves that declared policy. No restriction to only uniform-latent global content shapes is made, even though only those shapes occur in these particular bodies.

\begin{proposition}[Countable once-only scope boundary for infinite sharing]
Suppose a proposed infinite-execution extension has a countably supported once-only seed $L$ and two binary streams $X,Y\in\{0,1\}^{\mathbb N}$. Assume the entire streams are conditionally independent given $L$ and $X=Y$ almost surely. Then their common stream law is countably atomic. In particular it cannot be the iid fair-bit stream law, although every finite shared fair-bit prefix has a finite-seed representation.
\end{proposition}
\begin{proof}
For every seed value $\ell$ with positive mass, conditional equality holds almost surely, so the two conditional stream marginals agree. At coordinate $n$, let $p=\Pr(X_n=1\mid L=\ell)$. Conditional independence and equality give $p(1-p)=0$. Thus every coordinate is conditionally fixed. Since there are only countably many coordinates, the whole stream is almost surely one determined sequence $s_\ell$ under that seed value. The unconditional common law is therefore supported on the countable set of these sequences. The iid fair-bit law assigns each fixed sequence mass $\lim_n2^{-n}=0$, so cannot have countable support. A prefix of length $n$ can instead be sampled once from its $2^n$ possible bit strings and copied, proving the last claim.
\end{proof}
This is a boundary on a proposed infinite shared-execution interface, not a contradiction of the finite-terminating-tree theorem. Separate finite-prefix seed representations do not establish one countable once-only seed supporting the whole shared stream. Fresh jointly generated scopes or a different measurable latent interface would change the hypothesis and require their own native compatibility proof.

\section{A branching realization and finite diagnostics}
Take one recursive state $u$, one external input, and two global content shapes: a nullary shape $c_0$ and a binary shape $c_2$. At both the initial input and $u$, choose $c_0$ with probability $1/3$ and choose $c_2$ with outputs $(u,u)$ with probability $2/3$. The primitive source is finite and normalized. Its terminating raw witnesses are all finite full binary trees, not a unary encoding of their final outputs. The total finite-tree mass is the least solution of
\[
 q=\tfrac13+\tfrac23q^2,
\]
namely $q=1/2$, rather than the other solution $q=1$. There are infinitely many positive-weight witnesses, so this feedback does not lie in the old finite-raw normalized fragment. It lies in the explicitly constructed countable subprobability target.

Give every primitive context witness coefficient $k^2$ with augmentation $(x,y)\mapsto x$. Identity and the projection $(x,y)\mapsto(x,0)$ are different augmentation-preserving native endomorphisms with the same raw map and mass equation. Their tree-feedback maps are different: already the one-node terminal tree retains those different actual coefficient arrows, while larger trees retain their tensor products. Independently the two content shapes may carry the same augmented-vector coefficients in $\mathcal D=\mathbf{FinVect}_k/k$. The resulting global grade keeps the coefficients of all unused typed trees too; it is not determined by the output distribution. Every individual tree coefficient is finite-dimensional, although dimensions need not have a common bound and the whole family is countable.

The accompanying standard-library diagnostic enumerates all $1461$ coloured nullary/unary/binary tree shapes with one through four nodes, checking that cutting $U$ edges into $V$ blocks and flattening are mutually inverse on those shapes. It also computes seven exact rational height approximants for the displayed branching process and checks finite scalar aggregation when the terminal witness is split into weights $1/9$ and $2/9$. These are finite regressions, not proofs of the general categorical, countable or least-fixed-point statements. No TeX compilation or proof-assistant verification is claimed.

\section{What is not supplied by this construction}
The products of child effects in the tree formula are the original independent occurrence policy. Shared branch innovations require a separately retained joint conditional law; arbitrary component marginals do not determine it. The theorem does not claim that replacing the products by independently selected couplings preserves nesting. A viable shared extension must supply coupling data coherent under the very tree cuts used in the nesting proof, or retain a complete joint execution object rather than resample at each cut.

The scalar deficit $1-\sum_tp(t)$ can contain primitive missing mass, refusal, deadlock or infinite recursive execution. The finite-terminating-tree interface does not distinguish those causes and makes no claim that they are equal. A full semantic account must retain an appropriate failure/history interface or a separate infinite-tree measure. No infinite accumulated vector or content coefficient is constructed: every individual witness tree is finite, while the countable family of possible finite witnesses and its direct-sum augmentation are available in the declared target rig.

The positive increment is an independently constructed polynomial-content recursive-feedback operation, an original nonfaithful augmented-rig recovery, a native comparison theorem and a coherent full-grade nesting isomorphism. It is not a general solution for every original contentad, an instantaneous trace, an arbitrary shared-coupling completion, or a free universal completion of the whole relative construction. Those obligations remain explicit.
\section{Established mechanisms and remaining attribution limits}
The tree initial-algebra mechanism is established. Kock, \emph{Polynomial functors and trees}, Section 1.2, Theorem 1.2.5 and Proposition 1.2.8, constructs polynomial tree initial algebras and free polynomial monads by grafting. The present nonempty, internal-$U$-labelled collection and its block-cut isomorphism are explicit variants; no novelty claim is made for tree recursion or free polynomial structure. The additional work here is its action on complete native coefficient-decorated grades and its compatibility with the specified context/content crossing. Kock's ordinary polynomial theorem is not invoked as a proof of those actual coefficient or effect equations. Source checked: \url{https://arxiv.org/pdf/0807.2874}.

Likewise least probabilistic-polynomial fixed points are established branching-process machinery. Etessami, Stewart and Yannakakis, \emph{Polynomial Time Algorithms for Multi-Type Branching Processes and Stochastic Context-Free Grammars}, Introduction, formulates probabilistic polynomial systems with coefficient sums at most one and relates their least nonnegative solutions to termination/extinction probabilities. The present scalar height argument uses that familiar mechanism, while proving the countable raw-tree and native-arrow statements separately. No new algorithmic or classical branching-process claim is made. Source checked: \url{https://arxiv.org/pdf/1201.2374}.

These are targeted source checks, not an exhaustive novelty assessment. No claim is made that the entire coefficient-decorated mixed construction is absent from the literature.
\end{document}
